\chapter{Discusión}

Los resultados sostienen dos tipos de contribución. La contribución metodológica corresponde a la construcción y verificación del instrumento antes de utilizar sus mediciones para entrenar. La contribución empírica corresponde al comportamiento real de la plataforma y a la capacidad de transferencia de los clasificadores CPU y GPU entre familias algorítmicas.

Buena parte de la literatura revisada en el estado del arte (los métodos de selección de frecuencia óptima de Ali et al. \cite{Ali2023} y de clasificación consciente de DVFS de Guerreiro et al. \cite{Guerreiro2019}, o el marco de control por refuerzo de Wang et al. \cite{Wang2024}) asume, de forma razonable dado su alcance, que la escritura de frecuencia sobre el hardware subyacente funciona como el fabricante la documenta, y concentra su aporte en la capa de decisión: qué frecuencia elegir, o cómo aprenderla. Los resultados de este trabajo muestran que esa suposición debe verificarse empíricamente bajo carga real y no darse por sentada en ningún mecanismo del instrumento: en GPU, el bloqueo de reloj sostuvo dos objetivos distintos durante cientos de lecturas activas bajo carga; en CPU, la actuación se confirma correctamente solo cuando se restringen tanto los procesadores lógicos delegados como sus hermanos de \textit{SMT} (sección \ref{sec:resultados-cpu-dvfs}), un resultado que hubiera pasado por una lectura de límites exitosa sin serlo si solo se hubiera verificado la primera condición. La misma disciplina de no asumir una capacidad de hardware sin contrastarla rige el presupuesto de nueve contadores físicos simultáneos y no diez, ajustado por la reserva que hace el vigilante NMI del sistema operativo (sección \ref{sec:resultados-interferencia-uncore}), y la validación cruzada de los techos Roofline contra Intel Advisor, con una coincidencia de 12 a 15\% pese a partir de mecanismos de medición independientes (sección \ref{sec:resultados-ppico}). Generalizar esta disciplina de verificación directa, incluida una prueba de humo a escala reducida antes de cualquier campaña de producción, es un aporte metodológico de esta fase tan relevante como los datos concretos que produjo.

El resultado CPU confirma que la abundancia de intervalos no equivale a disponer de muchas unidades independientes de generalización. Los 1.31 millones de intervalos elegibles proceden de 30 familias algorítmicas y más de la mitad de ellas es homogénea respecto de la clase Roofline. La magnitud de esa distinción se pudo cuantificar con el mismo modelo y los mismos datos: repartir intervalos al azar produce 0.874 de exactitud balanceada, dejar fuera un kernel completo 0.733 y dejar fuera una familia completa 0.728 (sección \ref{sec:resultados-cpu-protocolo}). La primera cifra no mide clasificación de régimen sino reconocimiento de la carga, y es la que obtendría un protocolo que no separase por familia. Esta comparación es, en sí misma, un argumento metodológico a favor de reportar la validación por familia aun cuando produzca cifras notablemente menores.

El resultado también acota qué palancas quedan disponibles. El algoritmo de aprendizaje, el tamaño del vector de entrada y los hiperparámetros resultaron indiferentes dentro del intervalo de confianza, y un clasificador de vecinos más cercanos sin ajuste alcanza 0.684 con el mismo protocolo. La literatura revisada concentra su aporte en esa capa de decisión, de modo que conviene precisar el alcance del hallazgo: no contradice esos métodos, pero muestra que, en este catálogo y con estas señales, la frontera entre regímenes no se mueve al cambiar de modelo. Lo que sí mueve el resultado es la diversidad algorítmica del catálogo, con una pendiente estimada de 0.037 por duplicación del número de familias. Esa relación convierte la construcción del catálogo, y no la selección del modelo, en la actividad determinante del desempeño alcanzable.

La política selectiva ofrece una respuesta operativa acotada a esa incertidumbre, aunque su ganancia debe dimensionarse con cuidado. Con el umbral congelado en 0.85 el modelo decide sobre el 72.2\% de los intervalos y su exactitud balanceada por celda pasa de 0.728 a 0.738, una ganancia modesta por celda que se ve con más fuerza en las métricas agrupadas (la F1 macro pasa de 0.760 a 0.842). Abstenerse al azar en la misma proporción no cambia el resultado, de modo que la mejora proviene efectivamente de que la confianza señala casos difíciles. Aun así, la diferencia entre la F1 macro agrupada, de 0.842, y el desempeño de las familias individuales muestra que el buen resultado se concentra en familias transferibles y no constituye una garantía uniforme para una carga inédita.

Un resultado independiente del clasificador, pero decisivo para el objetivo de eficiencia, es que la clase Roofline por sí sola no determina la frecuencia de menor producto energía--retardo. En el procesador estudiado, unos 85~W de potencia de paquete no dependen del reloj de los núcleos, de modo que reducirlo de 3.2 a 0.8~GHz ahorra solo 15\% de potencia mientras alarga el tiempo en proporción a $f^{-\alpha}$. La política resultante es no actuar en ambas clases, y ni siquiera un selector perfecto que eligiera el mejor nivel por kernel supera un techo medio de 2.2\% de ganancia, con 22 de los 47 kernels por debajo del efecto mínimo relevante. Dos pruebas externas selladas, con una regla congelada antes de medir sobre familias nuevas ajenas al entrenamiento, confirmaron que no hay un subconjunto de la clase memory-bound sobre el que sea seguro actuar con la evidencia disponible (sección \ref{sec:resultados-cpu-potencia-diseno}). La utilidad del clasificador reside, por tanto, en decidir cuándo abstenerse de actuar, más que en repartir frecuencias por clase.

La abstención tampoco corrige los errores de alta confianza, porque la confianza del modelo está mal calibrada fuera de su familia de entrenamiento: declara 0.890 en promedio frente a una exactitud observada de 0.764, con un error de calibración esperado de 0.126. Entre los intervalos con confianza mayor o igual a 0.9, \textit{npb\_ft} acierta el 32\% y \textit{npb\_mg} el 35\%. Por ello la probabilidad estimada por XGBoost no debe interpretarse como una probabilidad calibrada de acierto para cualquier familia inédita; su función actual es delimitar una región de operación bajo el conjunto disponible. Para la salida \textit{revisar}, la política conservadora evita convertir la incertidumbre en una actuación DVFS potencialmente perjudicial; darle a esa espera una acción concreta, mediante una medición física breve de la carga en ejecución, es un diseño que corresponde al agente de la Fase 3 y no a esta.

El hallazgo de precisión mixta en el catálogo de GPU (sección \ref{sec:resultados-precision-gpu}) es relevante más allá de este trabajo concreto. La práctica de derivar la precisión de caracterización de un kernel a partir de su tipo de dato declarado en el código fuente, en lugar de la precisión efectivamente ejercitada por el hardware durante su ejecución, es razonable como primera aproximación, pero este trabajo encontró que puede diferir de lo medido en más del 70\% del trabajo aritmético de un kernel, y que esa diferencia puede situarse exactamente en el margen que decide una clasificación Roofline. Ningún trabajo de los revisados en el estado del arte reporta explícitamente esta verificación para GPU, lo que sugiere que la práctica de asumir la precisión declarada, sin contrastarla, podría no ser inusual en caracterizaciones similares.

El resultado sobre precisión de \texttt{gpu\_dgemm\_n4096} (sección \ref{sec:resultados-precision-gpu}) matiza, sin contradecir, el hallazgo de Guerreiro et al. sobre la sensibilidad de las cargas GPGPU a la ruta de ejecución elegida por bibliotecas de terceros \cite{Guerreiro2019}: una ruta acelerada por hardware especializado invalida efectivamente un método de caracterización diseñado para instrucciones aritméticas convencionales, pero no invalida necesariamente la magnitud medida, siempre que exista una vía de verificación independiente, en este caso el conteo analítico del propio algoritmo, contra la cual contrastarla.

El reanálisis histórico de GPU delimita el alcance de la clasificación ligera. El conjunto A se utilizó para seleccionar el candidato y el conjunto B, una campaña histórica independiente de agosto, se mantuvo para la prueba externa. La mejora frente al diagnóstico inicial no surgió de seleccionar retrospectivamente el mejor resultado de B, sino de excluir la familia miniBUDE con control de reloj no verificable, agregar las lecturas NVML por corrida, comparar representaciones robustas y ponderar cada familia de forma equivalente durante el entrenamiento. Las razones físicas derivadas, los percentiles, las medias recortadas y la dinámica temporal se evaluaron como alternativas. La representación estática basada en medianas e IQR fue la única que conservó utilidad fuera de su campaña de origen. El F1 macro externo de 0.663 supera ampliamente la regla que siempre predice la clase mayoritaria, cuyo F1 macro fue 0.426, aunque ambas tienen la misma exactitud de 0.743. La igualdad en exactitud no es contradictoria. El modelo recupera 18 corridas \textit{compute\_bound} que la regla mayoritaria nunca reconoce y, a cambio, confunde 18 corridas \textit{memory\_bound}.

La abstención calibrada ofrece una alternativa operativa coherente con esta limitación. La política no altera la etiqueta ni elimina retrospectivamente casos cercanos al \textit{ridge}. Solo utiliza la confianza calculada a partir de las señales NVML disponibles en producción y devuelve \textit{revisar} cuando esa confianza no supera el umbral seleccionado dentro del entrenamiento. Al retirar de la decisión automática la región en que ambas clases presentan señales agregadas similares, la F1 macro de las decisiones cubiertas aumenta a 0.734. Esta cifra se informa junto con la cobertura de 93.9\%, por lo que no se interpreta como una mejora de cobertura total ni como sustituto de la evaluación externa en B. Esta salida permite diferir una proporción pequeña de decisiones inciertas a una caracterización posterior, en lugar de forzar una clase con evidencia insuficiente.

La separabilidad parcial de la Figura \ref{fig:gpu-separabilidad-nvml} explica esta conducta. En los datos históricos, la etiqueta procede de una intensidad operacional por kernel obtenida fuera de línea, mientras que la entrada reúne las señales de todo el intervalo de la corrida. NVML describe actividad, potencia y estado operativo de toda la GPU, pero no observa directamente instrucciones ejecutadas, tráfico de memoria, ocupación, causas de espera o la identidad del kernel activo. Cuando una ejecución tiene subfases con comportamiento distinto, el agregado conserva su efecto medio y pierde el orden y los límites de esas subfases. Por ello, una familia \textit{compute\_bound} puede ocupar una región cercana a una familia \textit{memory\_bound}. La visualización t-SNE no se interpreta como prueba de generalización, dado que su objetivo es preservar vecindarios locales y sus distancias globales no son métricas físicas. La silueta por clase de 0.092 y el F1 macro de 0.616 obtenido por vecinos que excluyen la familia evaluada confirman el mismo diagnóstico. El conjunto contiene señal física parcial, pero no una frontera estable para todas las familias.

Esta limitación también precisa el papel de CUPTI. El perfilado de CUPTI o Nsight Compute puede producir una etiqueta Roofline por kernel en una campaña controlada, mientras que la traza de actividades aporta los tiempos de lanzamiento necesarios para alinear esa etiqueta con un segmento NVML. Esa combinación permite etiquetar una fase observada, pero no puede reconstruirse retrospectivamente cuando la campaña solo conserva muestras NVML y una etiqueta global de corrida. Estas capacidades fortalecen la construcción y la auditoría de la verdad de entrenamiento. No justifican incorporar contadores de perfilado al vector de producción, pues la repetición de kernels o las pasadas adicionales necesarias para recolectarlos contradicen el requisito de inferencia ligera. La continuación metodológica consiste en desbloquear la campaña, medir la intrusión de la traza CUPTI y repetir la evaluación externa con nuevas familias, manteniendo NVML como fuente de entrada del clasificador desplegable.
